
%% bare_jrnl.tex
%% V1.4b
%% 2015/08/26
%% by Michael Shell
%% see http://www.michaelshell.org/
%% for current contact information.
%%
%% This is a skeleton file demonstrating the use of IEEEtran.cls
%% (requires IEEEtran.cls version 1.8b or later) with an IEEE
%% journal paper.
%%
%% Support sites:
%% http://www.michaelshell.org/tex/ieeetran/
%% http://www.ctan.org/pkg/ieeetran
%% and
%% http://www.ieee.org/

%%*************************************************************************
%% Legal Notice:
%% This code is offered as-is without any warranty either expressed or
%% implied; without even the implied warranty of MERCHANTABILITY or
%% FITNESS FOR A PARTICULAR PURPOSE! 
%% User assumes all risk.
%% In no event shall the IEEE or any contributor to this code be liable for
%% any damages or losses, including, but not limited to, incidental,
%% consequential, or any other damages, resulting from the use or misuse
%% of any information contained here.
%%
%% All comments are the opinions of their respective authors and are not
%% necessarily endorsed by the IEEE.
%%
%% This work is distributed under the LaTeX Project Public License (LPPL)
%% ( http://www.latex-project.org/ ) version 1.3, and may be freely used,
%% distributed and modified. A copy of the LPPL, version 1.3, is included
%% in the base LaTeX documentation of all distributions of LaTeX released
%% 2003/12/01 or later.
%% Retain all contribution notices and credits.
%% ** Modified files should be clearly indicated as such, including  **
%% ** renaming them and changing author support contact information. **
%%*************************************************************************


% *** Authors should verify (and, if needed, correct) their LaTeX system  ***
% *** with the testflow diagnostic prior to trusting their LaTeX platform ***
% *** with production work. The IEEE's font choices and paper sizes can   ***
% *** trigger bugs that do not appear when using other class files.       ***                          ***
% The testflow support page is at:
% http://www.michaelshell.org/tex/testflow/



\documentclass[journal,onecolumn]{IEEEtran}
%
% If IEEEtran.cls has not been installed into the LaTeX system files,
% manually specify the path to it like:
% \documentclass[journal]{../sty/IEEEtran}





% Some very useful LaTeX packages include:
% (uncomment the ones you want to load)


% *** MISC UTILITY PACKAGES ***
%
%\usepackage{ifpdf}
% Heiko Oberdiek's ifpdf.sty is very useful if you need conditional
% compilation based on whether the output is pdf or dvi.
% usage:
% \ifpdf
%   % pdf code
% \else
%   % dvi code
% \fi
% The latest version of ifpdf.sty can be obtained from:
% http://www.ctan.org/pkg/ifpdf
% Also, note that IEEEtran.cls V1.7 and later provides a builtin
% \ifCLASSINFOpdf conditional that works the same way.
% When switching from latex to pdflatex and vice-versa, the compiler may
% have to be run twice to clear warning/error messages.






% *** CITATION PACKAGES ***
%
%\usepackage{cite}
% cite.sty was written by Donald Arseneau
% V1.6 and later of IEEEtran pre-defines the format of the cite.sty package
% \cite{} output to follow that of the IEEE. Loading the cite package will
% result in citation numbers being automatically sorted and properly
% "compressed/ranged". e.g., [1], [9], [2], [7], [5], [6] without using
% cite.sty will become [1], [2], [5]--[7], [9] using cite.sty. cite.sty's
% \cite will automatically add leading space, if needed. Use cite.sty's
% noadjust option (cite.sty V3.8 and later) if you want to turn this off
% such as if a citation ever needs to be enclosed in parenthesis.
% cite.sty is already installed on most LaTeX systems. Be sure and use
% version 5.0 (2009-03-20) and later if using hyperref.sty.
% The latest version can be obtained at:
% http://www.ctan.org/pkg/cite
% The documentation is contained in the cite.sty file itself.



\usepackage{graphicx}
%\usepackage{epstopdf}%to convert eps figures automatically NOT NEEDED ANYMORE
\usepackage{amsmath}
\usepackage{amsfonts}
\usepackage{amssymb}
\usepackage[amssymb]{SIunits}
%\usepackage{ifthen}
\usepackage{longtable}%flexible table definitions
%\usepackage{enumerate}%for easy definition of  number style

\usepackage{import}%to allow import of svg drawings
\usepackage{color}%to import svg drawings in color
\usepackage{algorithm,algpseudocode}% For algorithm environment
\usepackage{amsthm}
\usepackage{cite}
\usepackage{subcaption}
\usepackage{color}
\usepackage{enumitem}

% *** GRAPHICS RELATED PACKAGES ***
%
\ifCLASSINFOpdf
  % \usepackage[pdftex]{graphicx}
  % declare the path(s) where your graphic files are
  % \graphicspath{{../pdf/}{../jpeg/}}
  % and their extensions so you won't have to specify these with
  % every instance of \includegraphics
  % \DeclareGraphicsExtensions{.pdf,.jpeg,.png}
\else
  % or other class option (dvipsone, dvipdf, if not using dvips). graphicx
  % will default to the driver specified in the system graphics.cfg if no
  % driver is specified.
  % \usepackage[dvips]{graphicx}
  % declare the path(s) where your graphic files are
  % \graphicspath{{../eps/}}
  % and their extensions so you won't have to specify these with
  % every instance of \includegraphics
  % \DeclareGraphicsExtensions{.eps}
\fi
% graphicx was written by David Carlisle and Sebastian Rahtz. It is
% required if you want graphics, photos, etc. graphicx.sty is already
% installed on most LaTeX systems. The latest version and documentation
% can be obtained at: 
% http://www.ctan.org/pkg/graphicx
% Another good source of documentation is "Using Imported Graphics in
% LaTeX2e" by Keith Reckdahl which can be found at:
% http://www.ctan.org/pkg/epslatex
%
% latex, and pdflatex in dvi mode, support graphics in encapsulated
% postscript (.eps) format. pdflatex in pdf mode supports graphics
% in .pdf, .jpeg, .png and .mps (metapost) formats. Users should ensure
% that all non-photo figures use a vector format (.eps, .pdf, .mps) and
% not a bitmapped formats (.jpeg, .png). The IEEE frowns on bitmapped formats
% which can result in "jaggedy"/blurry rendering of lines and letters as
% well as large increases in file sizes.
%
% You can find documentation about the pdfTeX application at:
% http://www.tug.org/applications/pdftex





% *** MATH PACKAGES ***
%
%\usepackage{amsmath}
% A popular package from the American Mathematical Society that provides
% many useful and powerful commands for dealing with mathematics.
%
% Note that the amsmath package sets \interdisplaylinepenalty to 10000
% thus preventing page breaks from occurring within multiline equations. Use:
%\interdisplaylinepenalty=2500
% after loading amsmath to restore such page breaks as IEEEtran.cls normally
% does. amsmath.sty is already installed on most LaTeX systems. The latest
% version and documentation can be obtained at:
% http://www.ctan.org/pkg/amsmath





% *** SPECIALIZED LIST PACKAGES ***
%
%\usepackage{algorithmic}
% algorithmic.sty was written by Peter Williams and Rogerio Brito.
% This package provides an algorithmic environment fo describing algorithms.
% You can use the algorithmic environment in-text or within a figure
% environment to provide for a floating algorithm. Do NOT use the algorithm
% floating environment provided by algorithm.sty (by the same authors) or
% algorithm2e.sty (by Christophe Fiorio) as the IEEE does not use dedicated
% algorithm float types and packages that provide these will not provide
% correct IEEE style captions. The latest version and documentation of
% algorithmic.sty can be obtained at:
% http://www.ctan.org/pkg/algorithms
% Also of interest may be the (relatively newer and more customizable)
% algorithmicx.sty package by Szasz Janos:
% http://www.ctan.org/pkg/algorithmicx




% *** ALIGNMENT PACKAGES ***
%
%\usepackage{array}
% Frank Mittelbach's and David Carlisle's array.sty patches and improves
% the standard LaTeX2e array and tabular environments to provide better
% appearance and additional user controls. As the default LaTeX2e table
% generation code is lacking to the point of almost being broken with
% respect to the quality of the end results, all users are strongly
% advised to use an enhanced (at the very least that provided by array.sty)
% set of table tools. array.sty is already installed on most systems. The
% latest version and documentation can be obtained at:
% http://www.ctan.org/pkg/array


% IEEEtran contains the IEEEeqnarray family of commands that can be used to
% generate multiline equations as well as matrices, tables, etc., of high
% quality.




% *** SUBFIGURE PACKAGES ***
%\ifCLASSOPTIONcompsoc
%  \usepackage[caption=false,font=normalsize,labelfont=sf,textfont=sf]{subfig}
%\else
%  \usepackage[caption=false,font=footnotesize]{subfig}
%\fi
% subfig.sty, written by Steven Douglas Cochran, is the modern replacement
% for subfigure.sty, the latter of which is no longer maintained and is
% incompatible with some LaTeX packages including fixltx2e. However,
% subfig.sty requires and automatically loads Axel Sommerfeldt's caption.sty
% which will override IEEEtran.cls' handling of captions and this will result
% in non-IEEE style figure/table captions. To prevent this problem, be sure
% and invoke subfig.sty's "caption=false" package option (available since
% subfig.sty version 1.3, 2005/06/28) as this is will preserve IEEEtran.cls
% handling of captions.
% Note that the Computer Society format requires a larger sans serif font
% than the serif footnote size font used in traditional IEEE formatting
% and thus the need to invoke different subfig.sty package options depending
% on whether compsoc mode has been enabled.
%
% The latest version and documentation of subfig.sty can be obtained at:
% http://www.ctan.org/pkg/subfig




% *** FLOAT PACKAGES ***
%
%\usepackage{fixltx2e}
% fixltx2e, the successor to the earlier fix2col.sty, was written by
% Frank Mittelbach and David Carlisle. This package corrects a few problems
% in the LaTeX2e kernel, the most notable of which is that in current
% LaTeX2e releases, the ordering of single and double column floats is not
% guaranteed to be preserved. Thus, an unpatched LaTeX2e can allow a
% single column figure to be placed prior to an earlier double column
% figure.
% Be aware that LaTeX2e kernels dated 2015 and later have fixltx2e.sty's
% corrections already built into the system in which case a warning will
% be issued if an attempt is made to load fixltx2e.sty as it is no longer
% needed.
% The latest version and documentation can be found at:
% http://www.ctan.org/pkg/fixltx2e


%\usepackage{stfloats}
% stfloats.sty was written by Sigitas Tolusis. This package gives LaTeX2e
% the ability to do double column floats at the bottom of the page as well
% as the top. (e.g., "\begin{figure*}[!b]" is not normally possible in
% LaTeX2e). It also provides a command:
%\fnbelowfloat
% to enable the placement of footnotes below bottom floats (the standard
% LaTeX2e kernel puts them above bottom floats). This is an invasive package
% which rewrites many portions of the LaTeX2e float routines. It may not work
% with other packages that modify the LaTeX2e float routines. The latest
% version and documentation can be obtained at:
% http://www.ctan.org/pkg/stfloats
% Do not use the stfloats baselinefloat ability as the IEEE does not allow
% \baselineskip to stretch. Authors submitting work to the IEEE should note
% that the IEEE rarely uses double column equations and that authors should try
% to avoid such use. Do not be tempted to use the cuted.sty or midfloat.sty
% packages (also by Sigitas Tolusis) as the IEEE does not format its papers in
% such ways.
% Do not attempt to use stfloats with fixltx2e as they are incompatible.
% Instead, use Morten Hogholm'a dblfloatfix which combines the features
% of both fixltx2e and stfloats:
%
% \usepackage{dblfloatfix}
% The latest version can be found at:
% http://www.ctan.org/pkg/dblfloatfix




%\ifCLASSOPTIONcaptionsoff
%  \usepackage[nomarkers]{endfloat}
% \let\MYoriglatexcaption\caption
% \renewcommand{\caption}[2][\relax]{\MYoriglatexcaption[#2]{#2}}
%\fi
% endfloat.sty was written by James Darrell McCauley, Jeff Goldberg and 
% Axel Sommerfeldt. This package may be useful when used in conjunction with 
% IEEEtran.cls'  captionsoff option. Some IEEE journals/societies require that
% submissions have lists of figures/tables at the end of the paper and that
% figures/tables without any captions are placed on a page by themselves at
% the end of the document. If needed, the draftcls IEEEtran class option or
% \CLASSINPUTbaselinestretch interface can be used to increase the line
% spacing as well. Be sure and use the nomarkers option of endfloat to
% prevent endfloat from "marking" where the figures would have been placed
% in the text. The two hack lines of code above are a slight modification of
% that suggested by in the endfloat docs (section 8.4.1) to ensure that
% the full captions always appear in the list of figures/tables - even if
% the user used the short optional argument of \caption[]{}.
% IEEE papers do not typically make use of \caption[]'s optional argument,
% so this should not be an issue. A similar trick can be used to disable
% captions of packages such as subfig.sty that lack options to turn off
% the subcaptions:
% For subfig.sty:
% \let\MYorigsubfloat\subfloat
% \renewcommand{\subfloat}[2][\relax]{\MYorigsubfloat[]{#2}}
% However, the above trick will not work if both optional arguments of
% the \subfloat command are used. Furthermore, there needs to be a
% description of each subfigure *somewhere* and endfloat does not add
% subfigure captions to its list of figures. Thus, the best approach is to
% avoid the use of subfigure captions (many IEEE journals avoid them anyway)
% and instead reference/explain all the subfigures within the main caption.
% The latest version of endfloat.sty and its documentation can obtained at:
% http://www.ctan.org/pkg/endfloat
%
% The IEEEtran \ifCLASSOPTIONcaptionsoff conditional can also be used
% later in the document, say, to conditionally put the References on a 
% page by themselves.




% *** PDF, URL AND HYPERLINK PACKAGES ***
%
%\usepackage{url}
% url.sty was written by Donald Arseneau. It provides better support for
% handling and breaking URLs. url.sty is already installed on most LaTeX
% systems. The latest version and documentation can be obtained at:
% http://www.ctan.org/pkg/url
% Basically, \url{my_url_here}.




% *** Do not adjust lengths that control margins, column widths, etc. ***
% *** Do not use packages that alter fonts (such as pslatex).         ***
% There should be no need to do such things with IEEEtran.cls V1.6 and later.
% (Unless specifically asked to do so by the journal or conference you plan
% to submit to, of course. )

\usepackage{graphicx}
%\usepackage{epstopdf}%to convert eps figures automatically NOT NEEDED ANYMORE
\usepackage{amsmath}
\usepackage{amsfonts}
\usepackage{amssymb}
\usepackage[amssymb]{SIunits}
%\usepackage{ifthen}
\usepackage{longtable}%flexible table definitions
%\usepackage{enumerate}%for easy definition of  number style

\usepackage{import}%to allow import of svg drawings
\usepackage{color}%to import svg drawings in color
\usepackage{algorithm,algpseudocode}% For algorithm environment
\usepackage{amsthm}
\usepackage{commath}
\usepackage{subcaption} % not 'subfigure'
\usepackage{lipsum}% http://ctan.org/pkg/lipsum
\usepackage{multicol}% http://ctan.org/pkg/multicols
\usepackage{multirow}

%\newcolumntype{C}[1]{>{\centering}m{#1}}
\newcommand{\sminus}{\,{-}\,}
\newcommand{\splus}{\,{+}\,}
\newcommand{\sequal}{\,{=}\,}

% correct bad hyphenation here
\hyphenation{op-tical net-works semi-conduc-tor}


% Definition of environments
\newtheorem{mydef}{Definition}
\newtheorem{myremark}{Remark}
\newtheorem{mytheo}{Theorem}
\newtheorem{mycorollary}{Corollary}
\newtheorem{myprop}{Proposition}
\newtheorem{myproperty}{Property}
\newtheorem{myassump}{Assumption}


% Bold Symbols
\newcommand{\eqnarrayeq}{\!\!\!=\!\!\!}
\newcommand{\lk}{ \left\{ }
\newcommand{\rk}{ \right\} }
\newcommand{\la}{ \left\langle }
\newcommand{\ra}{ \right\rangle }
\newcommand{\tr}{ {\mbox{trace}} }
\newcommand{\rond}[1]{\frac{\partial}{\partial #1}}
\newcommand{\varrond}[2]{\frac{\partial #1}{\partial #2}}
\newcommand{\mb}{\left(}   % Matrix Begin
\newcommand{\me}{\right)}  % Matrix End
\newcommand{\argmin}{\mathop{\mbox{\rm argmin}}}
\newcommand{\argmax}{\mathop{\engmbox{\rm argmax}}}
\newcommand{\lmat}{\left[} % Left matrix
\newcommand{\rmat}{\right]}% Right matrix
\newcommand{\myrightarrow}{\mathop{\longrightarrow}}
\newcommand{\spl}{\mbox{sp}}


\newcommand{\Rbb}{{\mathbb{R}}}
\newcommand{\Zbb}{{\mathbb{Z}}}
\newcommand{\Ab}{{\bf A}}
\newcommand{\Tb}{{\bf T}}
\newcommand{\Hb}{{\bf H}}
\newcommand{\Cb}{{\bf C}}
\newcommand{\Rb}{{\bf R}}
\newcommand{\cb}{{\bf c}}
\newcommand{\wb}{{\bf w}}
\newcommand{\qb}{{\bf q}}
\newcommand{\Qb}{{\bf Q}}
\newcommand{\pb}{{\bf p}}
\newcommand{\ob}{{\bf o}}
\newcommand{\db}{{\bf d}}
\newcommand{\sbb}{{\bf s}}
\newcommand{\ab}{{\bf a}}
\newcommand{\xb}{{\bf x}}
\newcommand{\rb}{{\bf r}}
\newcommand{\yb}{{\bf y}}
\newcommand{\zb}{{\bf z}}
\newcommand{\Zb}{{\bf Z}}
\newcommand{\eb}{{\bf e}}
\newcommand{\nb}{{\bf n}}
\newcommand{\Db}{{\bf D}}
\newcommand{\Bb}{{\bf B}}
\newcommand{\Gb}{{\bf G}}
\newcommand{\Ob}{{\bf O}}
\newcommand{\Xb}{{\bf X}}
\newcommand{\Yb}{{\bf Y}}
\newcommand{\Eb}{{\bf E}}
\newcommand{\Wb}{{\bf W}}
\newcommand{\bze}{\boldmath 0}
\newcommand{\soft}[2]{{\cal S}_{#2}({#1})}
\newcommand{\nz}[1]{\|#1\|_0} % Norm Zero
\newcommand{\no}[1]{\|#1\|_1} % Norm One
\newcommand{\nf}[1]{\|#1\|_F} % Norm Fro
\newcommand{\nt}[1]{\|#1\|_2}
\newcommand{\np}[1]{\|#1\|_p}
\newcommand{\nq}[1]{\|#1\|_q}
\newcommand{\ninf}[1]{\|#1\|_{\infty}}
\newcommand{\Dc}{{\cal D}}
\newcommand{\Sc}{{\cal S}}
\newcommand{\Xc}{{\cal X}}
\newcommand{\cc}{{\cal c}}
\newcommand{\sac}{{\cal SA}}
\newcommand{\ud}{\,\mathrm{d}}
\newcommand{\Tetab}{{\mbox{\boldmath $\Theta$}}}
\newcommand{\Sigb}{{\mbox{\boldmath $\Sigma$}}}
\newcommand{\phib}{{\mbox{\boldmath $\phi$}}}
\newcommand{\psib}{{\mbox{\boldmath $\psi$}}}
\newcommand{\Psib}{{\mbox{\boldmath $\Psi$}}}
\newcommand{\Phib}{{\mbox{\boldmath $\Phi$}}}
\newcommand{\sign}{\mbox{{sgn}}}
\newcommand{\updt}{\mbox{{Update}}}
\newcommand{\dB}{\mbox{{dB}}}
\newcommand{\Tr}{\mbox{{Tr}}}
\newcommand{\diag}{\mbox{{diag}}}
\newcommand{\supp}{\mbox{{supp}}}
\newcommand{\off}{\mbox{{off}}}
\newcommand{\Thr}{\mbox{{Thresh}}}
\newcommand{\argg}{\mbox{{arg}}}
\newcommand{\nmz}{\mbox{{normalize}}}
\newcommand{\sprk}{\mbox{{spark}}}
\newcommand{\snr}{\mbox{{SNR}}}
\newcommand{\isnr}{\mbox{{ISNR}}}
\newcommand{\psnr}{\mbox{{PSNR}}}
\newcommand{\sbt}{\mbox{{subject~to}}}
\newcommand{\wrt}{\mbox{{w.r.t.}}}
\newcommand{\sbts}{\mbox{{s.t.}}}
\newcommand{\Xpinv}{{\Xb^{\!\dag}}}
\newcommand{\Wpinv}{{\Wb^{\!\dag}}}
\newcommand{\Bbh}{{\mathbf{\hat B}}}
\newcommand{\ybh}{{\mathbf{\hat y}}}
\newcommand{\xbh}{{\mathbf{\hat x}}}
\newcommand{\Rbh}{{\mathbf{\hat x}}}
\newcommand{\Bbs}{{\mathbf{\scriptsize B}}}
\newcommand{\bompcal}{\mbox{{Batch-OMP}}}
\newcommand{\Ib}{{\bf I}}
\newcommand{\Ei}{{\mbox{$\cal E$}}}
\newcommand{\RR}{{\Bbb R}}
\newcommand{\bijk}{b_{ij}^{(k)}}
\newcommand{\mub}{{\mbox{\boldmath $\mu$}}}
\newcommand{\Lamb}{{\mbox{\boldmath $\Lambda$}}}
\newcommand{\Omegb}{{\mbox{\boldmath $\Omega$}}}
\newcommand{\betab}{{\mbox{\boldmath $\beta$}}}
\newcommand{\Deltab}{{\mbox{\boldmath $\Delta$}}}
\newcommand{\deltab}{{\mbox{\boldmath $\delta$}}}
\newcommand{\etab}{{\mbox{\boldmath $\eta$}}}
\newcommand{\alphab}{{\mbox{\boldmath $\alpha$}}}
\newcommand{\gammab}{{\mbox{\boldmath $\gamma$}}}
\newcommand{\Gamab}{{\mbox{\boldmath $\Gamma$}}}
\newcommand{\xib}{{\mbox{\boldmath $\xi$}}}
\newcommand{\Ub}{{\mathbf U}}
\newcommand{\ub}{{\mathbf u}}
\newcommand{\vb}{{\mathbf v}}
\newcommand{\Vb}{{\mathbf V}}
\newcommand{\Sb}{{\mathbf S}}

% Calligraphic symbols
\newcommand{\Cc}{{\cal C}}
\newcommand{\Hc}{{\cal H}}
\newcommand{\Ic}{{\cal I}}
\newcommand{\Ac}{{\cal A}}
\newcommand{\Nc}{{\cal N}}
\newcommand{\cL}{{\cal L}}
\newcommand{\cU}{{\cal U}}
\newcommand{\cE}{{\cal E}}
\newcommand{\cG}{{\cal G}}
\newcommand{\Pc}{{\cal P}}
% News Commands
\newsavebox\mybox
\newenvironment{myfbox}{%
\begin{lrbox}{\mybox}%
\begin{minipage}{\dimexpr(\textwidth-2\fboxsep-2\fboxrule)}
}{%
\end{minipage}
\end{lrbox}%
\vskip10pt
\noindent
\fbox{\usebox\mybox}%
\vskip10pt
}
\newenvironment{myshadowbox}{%
\begin{lrbox}{\mybox}%
\begin{minipage}{\dimexpr(\textwidth-2\fboxsep-2\fboxrule-\shadowsize)}
}{%
\end{minipage}
\end{lrbox}%
\vskip10pt
\noindent
\shadowbox{\usebox\mybox}%
\vskip10pt
}


\newcommand{\idx}[1]{\index{#1}#1}
%% می تواند برای ارایه نکات در محیط itemize به کار رود، روند این کار به این صورت است،  (شکل یک تیر)
\newcommand{\arcO}{\noindent\textcolor{red}{\Large\ding{247}}\;}
%% این شکل می‌تواند برای بیان مزایای یک قضیه بکار رود (شکل تیک)
\newcommand{\tickO}{\noindent\textcolor{green}{\Large\ding{52}}\;}
%% برای  بیان معایب و یا نکات منفی (شکل یک ضربدر)
\newcommand{\X}{\item[\Large\color{red}\ding{56}]}
\newcommand{\XO}{\noindent\textcolor{red}{\LARGE\ding{56}}\;}
%% بیان موارد یک قضیه (شکل یک دست)
\newcommand{\handO}{\noindent\textcolor{blue}{\LARGE\ding{45}}\;}
%% برای مواردی که: این موارد شامل .... می شود، توسط عناصر زیر مشخص می شود (شکل یک درخت)
\newcommand{\treeO}{\noindent\textcolor{ForestGreen}{\Large\ding{171}}\;}
%% برای این که چند مورد را تعریف کنیم (علامت دست که دو گرفته)
\newcommand{\twoO}{\noindent\textcolor{blue}{\LARGE\ding{44}}\;}
%% (شکل یک قیچی)
\newcommand{\sciO}{\noindent\textcolor{OrangeRed}{\footnotesize\ding{108}}\;}

\newcommand{\teleO}{\noindent\textcolor{Bittersweet}{\huge\ding{37}}\;}

\begin{document}
%
% paper title
% Titles are generally capitalized except for words such as a, an, and, as,
% at, but, by, for, in, nor, of, on, or, the, to and up, which are usually
% not capitalized unless they are the first or last word of the title.
% Linebreaks \\ can be used within to get better formatting as desired.
% Do not put math or special symbols in the title.

%\title{Progressive Learning Network \\
%Progressive Learning for Systematic Design of \\ a Large Neural Network}
%\title{Progressive Learning Network - Use of Convex Optimization for Large Neural Network Design}
\title{Variational Deep Learning of Multiple Sinusoids}

%
%
% author names and IEEE memberships
% note positions of commas and nonbreaking spaces ( ~ ) LaTeX will not break
% a structure at a ~ so this keeps an author's name from being broken across
% two lines.
% use \thanks{} to gain access to the first footnote area
% a separate \thanks must be used for each paragraph as LaTeX2e's \thanks
% was not built to handle multiple paragraphs
%

%\author{Saikat~Chatterjee,~\IEEEmembership{Member,~IEEE,}
%		Alireza M. Javid,
%        Mostafa Sadeghi, \\
%        Partha P. Mitra,
%        Mikael Skoglund,~\IEEEmembership{Senior member,~IEEE}
% 
%
%        % <-this % stops a space
%%\thanks{M. Shell was with the Department
%%of Electrical and Computer Engineering, Georgia Institute of Technology, Atlanta,
%%GA, 30332 USA e-mail: (see http://www.michaelshell.org/contact.html).}% <-this % stops a space
%%\thanks{J. Doe and J. Doe are with Anonymous University.}% <-this % stops a space
%%\thanks{Manuscript received April 19, 2005; revised August 26, 2015.}
%}

\author{Lei, Fredrik, Magnus, Saikat}        
        

% note the % following the last \IEEEmembership and also \thanks - 
% these prevent an unwanted space from occurring between the last author name
% and the end of the author line. i.e., if you had this:
% 
% \author{....lastname \thanks{...} \thanks{...} }
%                     ^------------^------------^----Do not want these spaces!
%
% a space would be appended to the last name and could cause every name on that
% line to be shifted left slightly. This is one of those "LaTeX things". For
% instance, "\textbf{A} \textbf{B}" will typeset as "A B" not "AB". To get
% "AB" then you have to do: "\textbf{A}\textbf{B}"
% \thanks is no different in this regard, so shield the last } of each \thanks
% that ends a line with a % and do not let a space in before the next \thanks.
% Spaces after \IEEEmembership other than the last one are OK (and needed) as
% you are supposed to have spaces between the names. For what it is worth,
% this is a minor point as most people would not even notice if the said evil
% space somehow managed to creep in.



% The paper headers
\markboth{}%
{Shell \MakeLowercase{\textit{et al.}}: Bare Demo of IEEEtran.cls for IEEE Journals}
% The only time the second header will appear is for the odd numbered pages
% after the title page when using the twoside option.
% 
% *** Note that you probably will NOT want to include the author's ***
% *** name in the headers of peer review papers.                   ***
% You can use \ifCLASSOPTIONpeerreview for conditional compilation here if
% you desire.




% If you want to put a publisher's ID mark on the page you can do it like
% this:
%\IEEEpubid{0000--0000/00\$00.00~\copyright~2015 IEEE}
% Remember, if you use this you must call \IEEEpubidadjcol in the second
% column for its text to clear the IEEEpubid mark.



% use for special paper notices
%\IEEEspecialpapernotice{(Invited Paper)}




% make the title area
\maketitle

% As a general rule, do not put math, special symbols or citations
% in the abstract or keywords.
\begin{abstract}
We use the typical signal model - sum of sinusoids, as used in standard MUSIC, ESPRIT type of algorithms. And then develop a deep neural network based system, when the number of measurement samples is large. The measurement samples can be non-uniform and lower than standard Nyquist rate.  
\end{abstract} 

% Note that keywords are not normally used for peerreview papers.
\begin{IEEEkeywords}

\end{IEEEkeywords}

\section{Usual Notations}

\textcolor{red}{Note that we have a complex valued signal model for the "sinusoids". Hence we should either work with complex data or make it real valued, say by concatenating real and imaginary parts. }

We use the standard relation of Gaussian log likelihood:
\begin{eqnarray}
\begin{array}{rl}
\log \mathcal{N}(x;\mu, \sigma^2) & = -1/2 \log (2\pi) -1/2 \log \sigma^2  - \frac{1}{2 \sigma^2} (x- \mu)^2
\end{array}
\end{eqnarray}


\section{Signal Model}

The standard sinusoidal signal model and the measurement model are:
\begin{eqnarray}
    x_t = \sum_{k=1}^K \alpha_k \, e^{j( \omega_k t + \phi_k)}, \,\,\, y_t = x_t + w_t, w_t \sim \mathcal{N}(0,\sigma^2), t=t_1, t_2, \hdots, t_l, \hdots.
\end{eqnarray}
In practice, we can assume that the measurement sequence is received as a long sequence $\bar{\mathbf{y}} = y_{t_1}, y_{t_2}, \hdots, y_{t_n}, \hdots, y_{t_N}$, where $N$ is large. The sequence $\mathbf{y}$ corresponds to the sequence $\bar{\mathbf{x}} = x_{t_1}, x_{t_2}, \hdots, x_{t_n}, \hdots, x_{t_N}$. Let the parameters of  sinusoids be denoted by $3K$-dimensional vector $\pmb{\psi} = [\psi_1, \psi_2, \hdots, \psi_{k'}, \hdots, \psi_{3K}]^{\top} = [\alpha_1, \omega_1, \phi_1, \alpha_2, \omega_2, \phi_2, \hdots, \alpha_k, \omega_k, \phi_k, \hdots, \alpha_K, \omega_K, \phi_K]^{\top}$. Note that, given $\pmb{\psi}$, the sequence $\bar{\mathbf{x}}$ is fully deterministic. So, when appropriate, we use a notation $\bar{\mathbf{x}} \triangleq \bar{\mathbf{x}}(\pmb{\psi})$, and similarly $x_t \triangleq x_t(\pmb{\psi})$. In the Bayesian approach, our task is to find the posterior $p(\pmb{\psi}|\bar{\mathbf{y}})$. We have the following assumptions: 
\begin{eqnarray}
    p(\pmb{\psi})= \prod_{k'}^{3K} p(\psi_{k'}).
\end{eqnarray}


Note the following relations:
\begin{eqnarray}
\begin{array}{rl}
     p(y_t | \pmb{\psi}) & \triangleq p(y_t | x_t( \pmb{\psi})) = \mathcal{N}(\sum_{k=1}^K \alpha_kb \, e^{j( \omega_k t + \phi_k)}, \sigma^2)  \\
     p(\bar{\mathbf{y}}|\pmb{\psi}) & \triangleq p(\bar{\mathbf{y}} | \bar{\mathbf{x}}(\pmb{\psi})) = \prod_{n=1}^{N} p(y_{t_n}| \pmb{\psi}) = \prod_{n=1}^{N} p(y_{t_n}| x_{t_n}( \pmb{\psi}))  \\
     \log p(\bar{\mathbf{y}}|\pmb{\psi}) & =\sum_{n=1}^{N} \log p(y_{t_n}| \pmb{\psi}) 
\end{array}     
\end{eqnarray}


\subsection{A priori knowledge about appropriate distributions}

The quantities which we are interested in are amplitude, angular velocity and phase of sinusoids. The appropriate distributions to characterize them are Gaussian distribution, Maxwell–Boltzmann distribution, von Mises distribution. So use these distributions appropriately to represent variational distributions. In case we have problem, we have to use Gaussian distribution as Maxwell–Boltzmann and von Mises distributions have similarity to Gaussian distribution.

The challenge is how to sample from such distributions and design DNNs that allow gradient flow.


\subsection{Variational learning}

We first try to find a lower bound of the logarithmic evidence, as follows:

\begin{eqnarray}
    \begin{array}{rl}
      \log p(\bar{\mathbf{y}})   & = \displaystyle\int_{\pmb{\psi}} q(\pmb{\psi} | \bar{\mathbf{y}})  \log p(\bar{\mathbf{y}}) \, d\pmb{\psi} \\
         & = \displaystyle\int_{\pmb{\psi}} q(\pmb{\psi} | \bar{\mathbf{y}})  \log \frac{p(\bar{\mathbf{y}},\pmb{\psi})}{p(\pmb{\psi}|\bar{\mathbf{y}})} \, d\pmb{\psi} \\
         & = \displaystyle\int_{\pmb{\psi}} q(\pmb{\psi} | \bar{\mathbf{y}})  \log \left( \frac{p(\bar{\mathbf{y}},\pmb{\psi})} {q(\pmb{\psi} | \bar{\mathbf{y}})} . \frac{q(\pmb{\psi} | \bar{\mathbf{y}})}{p(\pmb{\psi}|\bar{\mathbf{y}})} \right) \, d\pmb{\psi} \\
         & = \underbrace{\displaystyle\int_{\pmb{\psi}} q(\pmb{\psi} | \bar{\mathbf{y}})  \log  \frac{p(\bar{\mathbf{y}},\pmb{\psi})} {q(\pmb{\psi} | \bar{\mathbf{y}})}  \, d\pmb{\psi}}_{\mathcal{L}(\bar{\mathbf{y}})} + \underbrace{\displaystyle\int_{\pmb{\psi}} q(\pmb{\psi} | \bar{\mathbf{y}})  \log \frac{q(\pmb{\psi} | \bar{\mathbf{y}})}{p(\pmb{\psi}|\bar{\mathbf{y}})} \, d\pmb{\psi} }_{\mathrm{D}_{\mathrm{KL}}(q\|p)} \\ 
    \end{array}
\end{eqnarray}
Here, $\mathrm{D}_{\mathrm{KL}}(q\|p)$ is the Kullback–Leibler (KL) divergence between $q(\pmb{\psi} | \bar{\mathbf{y}})$ and $p(\pmb{\psi} | \bar{\mathbf{y}})$. As $\mathrm{D}_{\mathrm{KL}}(q\|p) \geq 0$,  $\mathcal{L}(\bar{\mathbf{y}})$ is a lower bound of the logarithmic evidence $\log p(\bar{\mathbf{y}})$. This is also known evidence lower bound (ELBO). Maximization of $\mathcal{L}(\bar{\mathbf{y}})$ leads to decrease in $\mathrm{D}_{\mathrm{KL}}(q\|p)$, and $q(\pmb{\psi} | \bar{\mathbf{y}})$ approaches  $p(\pmb{\psi} | \bar{\mathbf{y}})$.


\subsection{Variational deep learning}

We now develop a variational deep learning approach where the variational distribution $q(\pmb{\psi} | \bar{\mathbf{y}})$ is provided by a deep neural network. We can rearrange the data as short multiple sequences or an array/tensor as we like, and then use sequence deep learning methods like RNNs, or Transformers, or array (or image) processing method like CNNs to have the parameters of $q(\pmb{\psi} | \bar{\mathbf{y}})$. Let the parameters of the chosen DNN be denoted by $\pmb{\theta}$. Therefore, we write $q(\pmb{\psi} | \bar{\mathbf{y}}) \triangleq q(\pmb{\psi} | \bar{\mathbf{y}}; \pmb{\theta})$. For simplicity, we assume $q(\pmb{\psi} | \bar{\mathbf{y}}; \pmb{\theta}) = \prod_{k'=1}^{3K} q(\psi_{k'} | \bar{\mathbf{y}}; \pmb{\theta})$. Now, we have the following relation:

\begin{eqnarray}
    \begin{array}{rl}
     \mathcal{L}(\bar{\mathbf{y}})  
         & \triangleq  \displaystyle\int_{\pmb{\psi}} q(\pmb{\psi} | \bar{\mathbf{y}} ; \pmb{\theta})  \log  \frac{p(\bar{\mathbf{y}},\pmb{\psi})} {q(\pmb{\psi} | \bar{\mathbf{y}} ; \pmb{\theta})}  \, d\pmb{\psi} \\
     
         & = \displaystyle\int_{\pmb{\psi}} q(\pmb{\psi} | \bar{\mathbf{y}} ; \pmb{\theta})  \log  \frac{p(\bar{\mathbf{y}}|\pmb{\psi}) p(\pmb{\psi})} {q(\pmb{\psi} | \bar{\mathbf{y}} ; \pmb{\theta})}  \, d\pmb{\psi} \\
         & = \displaystyle\int_{\pmb{\psi}} q(\pmb{\psi} | \bar{\mathbf{y}} ; \pmb{\theta})  \log  p(\bar{\mathbf{y}}|\pmb{\psi}) \, d\pmb{\psi} + \displaystyle\int_{\pmb{\psi}} q(\pmb{\psi} | \bar{\mathbf{y}} ; \pmb{\theta})  \log  \frac{ p(\pmb{\psi})} {q(\pmb{\psi} | \bar{\mathbf{y}} ; \pmb{\theta})}  \, d\pmb{\psi} \\ 

         & = \displaystyle\int_{\pmb{\psi}} q(\pmb{\psi} | \bar{\mathbf{y}} ; \pmb{\theta})  \log  p(\bar{\mathbf{y}}|\pmb{\psi}) \, d\pmb{\psi} - \displaystyle\int_{\pmb{\psi}} q(\pmb{\psi} | \bar{\mathbf{y}} ; \pmb{\theta})  \log  \frac{ q(\pmb{\psi} | \bar{\mathbf{y}} ; \pmb{\theta}) } {p(\pmb{\psi})}  \, d\pmb{\psi} \\ 

         & = \displaystyle\int_{\pmb{\psi}} \prod_{k'=1}^{3K} q(\psi_{k'} | \bar{\mathbf{y}}; \pmb{\theta}) \log \prod_{n=1}^{N}p(y_{t_n}| \pmb{\psi}) \, d\pmb{\psi} - \displaystyle\int_{\pmb{\psi}} \prod_{k'=1}^{3K} q(\psi_{k'} | \bar{\mathbf{y}}; \pmb{\theta})  \log  \frac{ \prod_{k'=1}^{3K} q(\psi_{k'} | \bar{\mathbf{y}}; \pmb{\theta})  } {\prod_{k'}^{3K} p(\psi_{k'})}  \, d\pmb{\psi} \\

         & = \displaystyle\int_{\pmb{\psi}} \prod_{k'=1}^{3K} q(\psi_{k'} | \bar{\mathbf{y}}; \pmb{\theta})  \sum_{n=1}^{N} \log p(y_{t_n}| \pmb{\psi}) \, d\pmb{\psi} - \sum_{k'=1}^{3K} \mathrm{D}_{\mathrm{KL}}(q(\psi_{k'} | \bar{\mathbf{y}}; \pmb{\theta}) \| p(\psi_{k'})) \\

         & =  \displaystyle\sum_{n=1}^{N} \mathbb{E}_{q(\pmb{\psi} | \bar{\mathbf{y}} ; \pmb{\theta})} [\log p(y_{t_n}| \pmb{\psi})] - \sum_{k'=1}^{3K} \mathrm{D}_{\mathrm{KL}}(q(\psi_{k'} | \bar{\mathbf{y}}; \pmb{\theta}) \| p(\psi_{k'})) \\

         & \triangleq \mathcal{L}(\bar{\mathbf{y}}; \pmb{\theta}, \sigma^2)
                  
    \end{array}
\end{eqnarray}

Therefore, the learning approach is:
\begin{eqnarray}
    \arg\max_{\pmb{\theta}, \sigma^2} \mathcal{L}(\bar{\mathbf{y}}; \pmb{\theta}, \sigma^2)
\end{eqnarray}

\textcolor{blue}{Fredrik please check the following argument:}
As $N$ is large and $\sum_{k'=1}^{3K} \mathrm{D}_{\mathrm{KL}}(q(\psi_{k'} | \bar{\mathbf{y}}; \pmb{\theta}) \| p(\psi_{k'})) \geq 0$, at the time of maximization-based learning we can use an approximation as: 
\begin{eqnarray}
    \arg\max_{\pmb{\theta}, \sigma^2} \mathcal{L}(\bar{\mathbf{y}}; \pmb{\theta}, \sigma^2) \approx \arg\max_{\pmb{\theta}, \sigma^2} \displaystyle\sum_{n=1}^{N} \mathbb{E}_{q(\pmb{\psi} | \bar{\mathbf{y}} ; \pmb{\theta})} [\log p(y_{t_n}| \pmb{\psi})].
\end{eqnarray}
In the above case, we do not need to have the prior $p(\pmb{\psi})= \prod_{k'}^{3K} p(\psi_{k'})$, or can use a non-informative prior.


\subsection{Profile-LS formulation for global frequency inference}

The equations already introduced above imply the following physical assumptions. The factorization of $p(\pmb{\psi})$ means that the sinusoidal parameters are a priori independent. The factorization of $p(\bar{\mathbf{y}}|\pmb{\psi})$ over time means that, conditioned on the sinusoidal parameters, the measurement noise is independent and identically distributed across observation times. The factorization of $q(\pmb{\psi}|\bar{\mathbf{y}};\pmb{\theta})$ means that the variational posterior neglects posterior dependence among different latent parameters. The use of a Gaussian likelihood means that the reconstruction error is modeled as additive Gaussian measurement noise.

For the profile-LS formulation, amplitude and phase are combined into a complex coefficient,
\begin{eqnarray}
    c_k \triangleq \alpha_k e^{j\phi_k},
    \qquad
    f_k \triangleq \frac{\omega_k}{2\pi},
\end{eqnarray}
so that
\begin{eqnarray}
    x_t = \sum_{k=1}^{K} c_k e^{j2\pi f_k t}.
\end{eqnarray}
The following additional assumptions are required.

\begin{myassump}[Non-overlapping patches]
The long observation sequence is partitioned into $B$ non-overlapping patches,
\begin{eqnarray}
    \bar{\mathbf{y}} = \{\mathbf{y}_1,\mathbf{y}_2,\ldots,\mathbf{y}_B\},
    \qquad
    \mathbf{y}_b = [y_{b,1},y_{b,2},\ldots,y_{b,L_b}]^{\top} \in \mathbb{C}^{L_b}.
\end{eqnarray}
Physically, every observation is used exactly once, so the patch likelihoods do not duplicate measurement information.
\end{myassump}

\begin{myassump}[One shared global frequency vector]
All patches observe the same fixed frequency vector,
\begin{eqnarray}
    \mathbf{f} = [f_1,f_2,\ldots,f_K]^{\top}.
\end{eqnarray}
Physically, the modal frequencies remain constant over the complete sequence, although the data are processed patch by patch.
\end{myassump}

\begin{myassump}[Patch-specific complex amplitudes]
Each patch has its own complex-amplitude vector,
\begin{eqnarray}
    \mathbf{c}_b = [c_{b,1},c_{b,2},\ldots,c_{b,K}]^{\top} \in \mathbb{C}^{K},
\end{eqnarray}
and
\begin{eqnarray}
    \mathbf{y}_b = \mathbf{\Phi}_b(\mathbf{f})\mathbf{c}_b + \mathbf{w}_b,
    \qquad
    [\mathbf{\Phi}_b(\mathbf{f})]_{n,k}=e^{j2\pi f_k t_{b,n}}.
\end{eqnarray}
Physically, the frequencies are globally fixed, while amplitude and phase may change between patches.
\end{myassump}

\begin{myassump}[Proper complex Gaussian noise]
The measurement noise satisfies
\begin{eqnarray}
    w_{b,n} \overset{\mathrm{iid}}{\sim} \mathcal{CN}(0,\sigma_w^2),
    \qquad
    \mathbf{w}_b \sim \mathcal{CN}(\mathbf{0},\sigma_w^2\mathbf{I}_{L_b}).
\end{eqnarray}
Physically, the real and imaginary noise components are zero-mean Gaussian, circularly symmetric, and independent across all sampled times and patches.
\end{myassump}

\begin{myassump}[Conditional independence across patches]
Conditioned on the shared frequency vector and the patch-specific amplitudes,
\begin{eqnarray}
    p(\bar{\mathbf{y}}|\mathbf{f},\mathbf{c}_{1:B})
    = \prod_{b=1}^{B}p(\mathbf{y}_b|\mathbf{f},\mathbf{c}_b).
\end{eqnarray}
Physically, after conditioning on the unknown sinusoidal parameters, one patch contains no additional stochastic information about another patch.
\end{myassump}

\begin{myassump}[Ridge least-squares profiling of complex amplitudes]
For a given frequency vector, the complex amplitude in patch $b$ is replaced by
\begin{eqnarray}
\begin{array}{rl}
    \widehat{\mathbf{c}}_b(\mathbf{f})
    & \triangleq \displaystyle\arg\min_{\mathbf{c}_b}
    \left\{
    \|\mathbf{y}_b-\mathbf{\Phi}_b(\mathbf{f})\mathbf{c}_b\|_2^2
    +\lambda\|\mathbf{c}_b\|_2^2
    \right\} \\
    & = \left(\mathbf{\Phi}_b^{\mathrm{H}}(\mathbf{f})\mathbf{\Phi}_b(\mathbf{f})+\lambda\mathbf{I}\right)^{-1}
    \mathbf{\Phi}_b^{\mathrm{H}}(\mathbf{f})\mathbf{y}_b.
\end{array}
\end{eqnarray}
Physically, amplitudes are nuisance parameters solved analytically for every sampled frequency vector rather than represented by a variational posterior.
\end{myassump}

\begin{myassump}[Uniform frequency prior on component-specific search bands]
For $f_k\in[\ell_k,u_k]$,
\begin{eqnarray}
    p(f_k)=\frac{1}{u_k-\ell_k}\mathbf{1}_{[\ell_k,u_k]}(f_k),
\end{eqnarray}
and
\begin{eqnarray}
    p(\mathbf{f})=\prod_{k=1}^{K}p(f_k).
\end{eqnarray}
Physically, only the admissible search interval of each component is known a priori, and all frequencies inside that interval are treated as equally plausible.
\end{myassump}

\begin{myassump}[Factorized local truncated-Gaussian frequency posteriors]
For every patch $b$, the DNN provides
\begin{eqnarray}
    q_b(\mathbf{f}|\mathbf{y}_b;\pmb{\theta})
    = \prod_{k=1}^{K}q_b(f_k|\mathbf{y}_b;\pmb{\theta}),
\end{eqnarray}
where
\begin{eqnarray}
    q_b(f_k|\mathbf{y}_b;\pmb{\theta})
    =
    \frac{
    \varphi\!\left((f_k-\mu_{b,k})/\sigma_{b,k}\right)
    }{
    \sigma_{b,k} Z_{b,k}
    }
    \mathbf{1}_{[\ell_k,u_k]}(f_k),
\end{eqnarray}
\begin{eqnarray}
    Z_{b,k}
    =
    \Phi\!\left(\frac{u_k-\mu_{b,k}}{\sigma_{b,k}}\right)
    -
    \Phi\!\left(\frac{\ell_k-\mu_{b,k}}{\sigma_{b,k}}\right).
\end{eqnarray}
Physically, each patch provides independent evidence about the same global frequencies, while every frequency remains inside its known search band.
\end{myassump}

\begin{myassump}[Global posterior obtained by posterior fusion]
The global variational posterior is formed from the local patch posteriors according to
\begin{eqnarray}
    q_{\mathrm{G}}(\mathbf{f}|\bar{\mathbf{y}};\pmb{\theta})
    \propto
    p(\mathbf{f})^{1-B}
    \prod_{b=1}^{B}q_b(\mathbf{f}|\mathbf{y}_b;\pmb{\theta}).
\end{eqnarray}
Physically, all non-overlapping patches contribute evidence about one shared frequency vector, while the common prior is counted only once.
\end{myassump}

\subsubsection{Derivation of the global frequency posterior}

\begin{eqnarray}
\begin{array}{rl}
p(\mathbf{f}|\bar{\mathbf{y}})
& = p(\mathbf{f}|\mathbf{y}_1,\ldots,\mathbf{y}_B) \\
& = p(\mathbf{f})\displaystyle\frac{p(\mathbf{y}_1,\ldots,\mathbf{y}_B|\mathbf{f})}{p(\mathbf{y}_1,\ldots,\mathbf{y}_B)}  \\
& = p(\mathbf{f})\displaystyle\frac{\prod_{b=1}^{B}p(\mathbf{y}_b|\mathbf{f})}{p(\mathbf{y}_1,\ldots,\mathbf{y}_B)} \\
& = \displaystyle\frac{p(\mathbf{f})}{p(\mathbf{y}_1,\ldots,\mathbf{y}_B)}\displaystyle\prod_{b=1}^{B}
\frac{p(\mathbf{f}|\mathbf{y}_b)p(\mathbf{y}_b)}{p(\mathbf{f})} \\
& \propto p(\mathbf{f})^{1-B}\displaystyle\prod_{b=1}^{B}p(\mathbf{f}|\mathbf{y}_b).
\end{array}
\end{eqnarray}

\begin{eqnarray}
\begin{array}{rl}
q_{\mathrm{G}}(\mathbf{f}|\bar{\mathbf{y}};\pmb{\theta})
& \triangleq \displaystyle\prod_{k=1}^{K} 
q_{\mathrm{G},k}(f_k|\bar{\mathbf{y}};\pmb{\theta}) \\
& \propto \displaystyle\prod_{k=1}^{K}
\left[
p(f_k)^{1-B}
\prod_{b=1}^{B}q_b(f_k|\mathbf{y}_b;\pmb{\theta})
\right]\\
& \propto \displaystyle\prod_{k=1}^{K}
\left[\mathbf{1}_{[\ell_k,u_k]}(f_k)
\displaystyle\prod_{b=1}^{B}
\exp\!\left[-\frac{(f_k-\mu_{b,k})^2}{2\sigma_{b,k}^2}\right]\right] \\
& \propto \displaystyle\prod_{k=1}^{K}
\left[\mathbf{1}_{[\ell_k,u_k]}(f_k)
\exp\!\left[
-\frac{1}{2}
\sum_{b=1}^{B}
\frac{(f_k-\mu_{b,k})^2}{\sigma_{b,k}^2}
\right]\right] \\
& \propto \displaystyle\prod_{k=1}^{K}
\left[ \mathbf{1}_{[\ell_k,u_k]}(f_k)
\exp\!\left[
-\frac{(f_k-\mu_{\mathrm{G},k})^2}{2\sigma_{\mathrm{G},k}^2}
\right]\right].
\end{array}
\end{eqnarray}


\begin{eqnarray}
    \sigma_{\mathrm{G},k}^{2}
    =\left(\sum_{b=1}^{B}\sigma_{b,k}^{-2}\right)^{-1},
\
    \mu_{\mathrm{G},k}
    =\sigma_{\mathrm{G},k}^{2}
    \sum_{b=1}^{B}\frac{\mu_{b,k}}{\sigma_{b,k}^{2}}
\end{eqnarray}

\begin{eqnarray}
    a_{\mathrm{G},k}
    \triangleq\frac{\ell_k-\mu_{\mathrm{G},k}}{\sigma_{\mathrm{G},k}},
    \qquad
    b_{\mathrm{G},k}
    \triangleq\frac{u_k-\mu_{\mathrm{G},k}}{\sigma_{\mathrm{G},k}},
    \qquad
    Z_{\mathrm{G},k}
    \triangleq
    \Phi(b_{\mathrm{G},k})-\Phi(a_{\mathrm{G},k}),
\end{eqnarray}

\begin{eqnarray}
    q_{\mathrm{G},k}(f_k|\bar{\mathbf{y}};\pmb{\theta})
    \propto
    \frac{
    \varphi\!\left((f_k-\mu_{\mathrm{G},k})/\sigma_{\mathrm{G},k}\right)
    }{
    \sigma_{\mathrm{G},k}Z_{\mathrm{G},k}
    }
    \mathbf{1}_{[\ell_k,u_k]}(f_k),
\end{eqnarray}


\begin{myassump}[Monte Carlo evaluation]
The likelihood expectation is approximated using $S$ samples,
\begin{eqnarray}
    \mathbf{f}^{(s)} \sim q_{\mathrm{G}}(\mathbf{f}|\bar{\mathbf{y}};\pmb{\theta}),
    \qquad s=1,2,\ldots,S.
\end{eqnarray}
Physically, every sampled global frequency vector is applied to all patches, while the patch-specific amplitudes are re-estimated for that sample.
\end{myassump}

The replacement of the amplitude posterior by $\widehat{\mathbf{c}}_b(\mathbf{f})$ changes the original joint frequency--amplitude ELBO into a frequency-only profile objective. The derivation below therefore continues the original ELBO after analytically profiling out the nuisance amplitudes.

\subsection{Continuation of the ELBO derivation for profile LS}

\begin{eqnarray}
\begin{array}{rl} p\!\left(\mathbf{y}_b|\mathbf{f},\widehat{\mathbf{c}}_b(\mathbf{f})\right) 
& \triangleq
\widetilde{p}(\mathbf{y}_b|\mathbf{f}) 
\end{array}
\end{eqnarray}

\begin{eqnarray}
\begin{array}{rl}
\log \widetilde{p}(\bar{\mathbf{y}}|\mathbf{f})
& = \displaystyle\sum_{b=1}^{B}\log \widetilde{p}(\mathbf{y}_b|\mathbf{f}) 
= -\displaystyle\sum_{b=1}^{B}L_b\log(\pi\sigma_w^2)
-\frac{1}{\sigma_w^2}
\displaystyle\sum_{b=1}^{B}
\left\|\mathbf{y}_b-\mathbf{\Phi}_b(\mathbf{f})\widehat{\mathbf{c}}_b(\mathbf{f})\right\|_2^2.
\end{array}
\end{eqnarray}

\begin{eqnarray}
\begin{array}{rl}
\mathcal{L}_{\mathrm{profile}}(\bar{\mathbf{y}};\pmb{\theta})
& \triangleq \displaystyle\int_{\mathbf{f}}
q_{\mathrm{G}}(\mathbf{f}|\bar{\mathbf{y}};\pmb{\theta})
\log\frac{\widetilde{p}(\bar{\mathbf{y}}|\mathbf{f})p(\mathbf{f})}
{q_{\mathrm{G}}(\mathbf{f}|\bar{\mathbf{y}};\pmb{\theta})}
\,d\mathbf{f} \\
& = \displaystyle\int_{\mathbf{f}}
q_{\mathrm{G}}(\mathbf{f}|\bar{\mathbf{y}};\pmb{\theta})
\log\widetilde{p}(\bar{\mathbf{y}}|\mathbf{f})\,d\mathbf{f}
-\displaystyle\int_{\mathbf{f}}
q_{\mathrm{G}}(\mathbf{f}|\bar{\mathbf{y}};\pmb{\theta})
\log\frac{q_{\mathrm{G}}(\mathbf{f}|\bar{\mathbf{y}};\pmb{\theta})}{p(\mathbf{f})}\,d\mathbf{f} \\
& = \mathbb{E}_{q_{\mathrm{G}}(\mathbf{f}|\bar{\mathbf{y}};\pmb{\theta})}
\left[\log\widetilde{p}(\bar{\mathbf{y}}|\mathbf{f})\right]
-\mathrm{D}_{\mathrm{KL}}\!\left(
q_{\mathrm{G}}(\mathbf{f}|\bar{\mathbf{y}};\pmb{\theta})\|p(\mathbf{f})
\right) \\
& = \mathbb{E}_{q_{\mathrm{G}}(\mathbf{f}|\bar{\mathbf{y}};\pmb{\theta})}
\left[
-\displaystyle\sum_{b=1}^{B}L_b\log(\pi\sigma_w^2)
-\frac{1}{\sigma_w^2}
\displaystyle\sum_{b=1}^{B}
\left\|\mathbf{y}_b-\mathbf{\Phi}_b(\mathbf{f})\widehat{\mathbf{c}}_b(\mathbf{f})\right\|_2^2
\right] -\mathrm{D}_{\mathrm{KL}}\!\left(
q_{\mathrm{G}}(\mathbf{f}|\bar{\mathbf{y}};\pmb{\theta})\|p(\mathbf{f})
\right) \\
& = \displaystyle-\sum_{b=1}^{B}L_b\log(\pi\sigma_w^2)
-\frac{1}{\sigma_w^2}
\mathbb{E}_{q_{\mathrm{G}}(\mathbf{f}|\bar{\mathbf{y}};\pmb{\theta})}
\left[
\displaystyle\sum_{b=1}^{B}
\left\|\mathbf{y}_b-\mathbf{\Phi}_b(\mathbf{f})\widehat{\mathbf{c}}_b(\mathbf{f})\right\|_2^2
\right] -\mathrm{D}_{\mathrm{KL}}\!\left(
q_{\mathrm{G}}(\mathbf{f}|\bar{\mathbf{y}};\pmb{\theta})\|p(\mathbf{f})
\right) \\
& \approx \displaystyle-\sum_{b=1}^{B}L_b\log(\pi\sigma_w^2) -\displaystyle\frac{1}{S\sigma_w^2}
\sum_{s=1}^{S}\sum_{b=1}^{B}
\left\|
\mathbf{y}_b-\mathbf{\Phi}_b(\mathbf{f}^{(s)})
\widehat{\mathbf{c}}_b(\mathbf{f}^{(s)})
\right\|_2^2 -\mathrm{D}_{\mathrm{KL}}\!\left(
q_{\mathrm{G}}(\mathbf{f}|\bar{\mathbf{y}};\pmb{\theta})\|p(\mathbf{f})
\right) \\
& = \displaystyle-\sum_{b=1}^{B}L_b\log(\pi\sigma_w^2) -\displaystyle\frac{1}{S\sigma_w^2}
\sum_{s=1}^{S}\sum_{b=1}^{B}
\left\|
\mathbf{y}_b-\mathbf{\Phi}_b(\mathbf{f}^{(s)})
\left(
\mathbf{\Phi}_b^{\mathrm{H}}(\mathbf{f}^{(s)})
\mathbf{\Phi}_b(\mathbf{f}^{(s)})+\lambda\mathbf{I}
\right)^{-1}
\mathbf{\Phi}_b^{\mathrm{H}}(\mathbf{f}^{(s)})\mathbf{y}_b
\right\|_2^2 \\
& \quad -\mathrm{D}_{\mathrm{KL}}\!\left(
q_{\mathrm{G}}(\mathbf{f}|\bar{\mathbf{y}};\pmb{\theta})\|p(\mathbf{f})
\right).
\end{array}
\end{eqnarray}

\subsubsection{Expansion of the global frequency KL divergence}

\begin{eqnarray}
\begin{array}{rl}
\mathrm{D}_{\mathrm{KL}}\!\left(
q_{\mathrm{G}}(\mathbf{f}|\bar{\mathbf{y}};\pmb{\theta})
\|p(\mathbf{f})
\right)
& = \displaystyle\sum_{k=1}^{K}
\mathrm{D}_{\mathrm{KL}}\!\left(
q_{\mathrm{G},k}(f_k|\bar{\mathbf{y}};\pmb{\theta})
\|p(f_k)
\right).
\end{array}
\end{eqnarray}

\begin{eqnarray}
\begin{array}{rl}
\mathrm{D}_{\mathrm{KL}}\!\left(
q_{\mathrm{G},k}\|p(f_k)
\right)
& = \displaystyle\int_{\ell_k}^{u_k}
q_{\mathrm{G},k}(f_k)
\log\frac{q_{\mathrm{G},k}(f_k)}{p(f_k)}\,df_k \\
& = \displaystyle\int_{\ell_k}^{u_k}
q_{\mathrm{G},k}(f_k)
\log\left[
\frac{u_k-\ell_k}{\sigma_{\mathrm{G},k}Z_{\mathrm{G},k}}
\varphi\!\left(
\frac{f_k-\mu_{\mathrm{G},k}}{\sigma_{\mathrm{G},k}}
\right)
\right]df_k \\
& = \log\frac{u_k-\ell_k}
{\sigma_{\mathrm{G},k}Z_{\mathrm{G},k}\sqrt{2\pi}}
-\frac{1}{2}
\mathbb{E}_{q_{\mathrm{G},k}}
\left[
\left(
\frac{f_k-\mu_{\mathrm{G},k}}{\sigma_{\mathrm{G},k}}
\right)^2
\right].
\end{array}
\end{eqnarray}

\begin{eqnarray}
\mathbb{E}_{q_{\mathrm{G},k}}
\left[
\left(
\frac{f_k-\mu_{\mathrm{G},k}}{\sigma_{\mathrm{G},k}}
\right)^2
\right]
=
1+
\frac{
a_{\mathrm{G},k}\varphi(a_{\mathrm{G},k})
-b_{\mathrm{G},k}\varphi(b_{\mathrm{G},k})
}{Z_{\mathrm{G},k}},
\end{eqnarray}

\begin{eqnarray}
\begin{array}{rl}
\mathrm{D}_{\mathrm{KL}}\!\left(
q_{\mathrm{G},k}\|p(f_k)
\right)
& = \log\frac{u_k-\ell_k}
{\sigma_{\mathrm{G},k}Z_{\mathrm{G},k}\sqrt{2\pi}}
-\frac{1}{2}
\left[
1+
\frac{
a_{\mathrm{G},k}\varphi(a_{\mathrm{G},k})
-b_{\mathrm{G},k}\varphi(b_{\mathrm{G},k})
}{Z_{\mathrm{G},k}}
\right] \\
& = \log\frac{u_k-\ell_k}
{\sigma_{\mathrm{G},k}Z_{\mathrm{G},k}\sqrt{2\pi e}}
-
\frac{
a_{\mathrm{G},k}\varphi(a_{\mathrm{G},k})
-b_{\mathrm{G},k}\varphi(b_{\mathrm{G},k})
}{2Z_{\mathrm{G},k}}.
\end{array}
\end{eqnarray}

\begin{eqnarray}
\begin{array}{rl}
\mathrm{D}_{\mathrm{KL}}\!\left(
q_{\mathrm{G}}(\mathbf{f}|\bar{\mathbf{y}};\pmb{\theta})
\|p(\mathbf{f})
\right)
& = \displaystyle\sum_{k=1}^{K}
\left[
\log\frac{u_k-\ell_k}
{\sigma_{\mathrm{G},k}Z_{\mathrm{G},k}\sqrt{2\pi e}}
-
\frac{
a_{\mathrm{G},k}\varphi(a_{\mathrm{G},k})
-b_{\mathrm{G},k}\varphi(b_{\mathrm{G},k})
}{2Z_{\mathrm{G},k}}
\right].
\end{array}
\end{eqnarray}






%\appendices
%\section{Proof of the First Zonklar Equation}
%Appendix one text goes here.

% you can choose not to have a title for an appendix
% if you want by leaving the argument blank


% use section* for acknowledgment
%\section*{Acknowledgment}
%
%
%The authors would like to thank...


% Can use something like this to put references on a page
% by themselves when using endfloat and the captionsoff option.
\ifCLASSOPTIONcaptionsoff
  \newpage
\fi



% trigger a \newpage just before the given reference
% number - used to balance the columns on the last page
% adjust value as needed - may need to be readjusted if
% the document is modified later
%\IEEEtriggeratref{8}
% The "triggered" command can be changed if desired:
%\IEEEtriggercmd{\enlargethispage{-5in}}

% references section

% can use a bibliography generated by BibTeX as a .bbl file
% BibTeX documentation can be easily obtained at:
% http://mirror.ctan.org/biblio/bibtex/contrib/doc/
% The IEEEtran BibTeX style support page is at:
% http://www.michaelshell.org/tex/ieeetran/bibtex/
%\bibliographystyle{IEEEtran}
% argument is your BibTeX string definitions and bibliography database(s)
%\bibliography{IEEEabrv,../bib/paper}
%
% <OR> manually copy in the resultant .bbl file
% set second argument of \begin to the number of references
% (used to reserve space for the reference number labels box)

\bibliographystyle{IEEEbib}
%\bibliography{ref,ref_alireza,biblio_saikat_ANN}
% biography section
% 
% If you have an EPS/PDF photo (graphicx package needed) extra braces are
% needed around the contents of the optional argument to biography to prevent
% the LaTeX parser from getting confused when it sees the complicated
% \includegraphics command within an optional argument. (You could create
% your own custom macro containing the \includegraphics command to make things
% simpler here.)
%\begin{IEEEbiography}[{\includegraphics[width=1in,height=1.25in,clip,keepaspectratio]{mshell}}]{Michael Shell}
% or if you just want to reserve a space for a photo:

%\begin{IEEEbiography}{Michael Shell}
%Biography text here.
%\end{IEEEbiography}
%
%% if you will not have a photo at all:
%\begin{IEEEbiographynophoto}{John Doe}
%Biography text here.
%\end{IEEEbiographynophoto}

% insert where needed to balance the two columns on the last page with
% biographies
%\newpage

%\begin{IEEEbiographynophoto}{Jane Doe}
%Biography text here.
%\end{IEEEbiographynophoto}
%
% You can push biographies down or up by placing
% a \vfill before or after them. The appropriate
% use of \vfill depends on what kind of text is
% on the last page and whether or not the columns
% are being equalized.

%\vfill

% Can be used to pull up biographies so that the bottom of the last one
% is flush with the other column.
%\enlargethispage{-5in}



% that's all folks
\end{document}


